From Kepler's 3\textsuperscript{rd} Law, $T^2 = a^3$ in units of Earth years and au respectively.
\[
\therefore T_{\mathrm{Mercury}} = \sqrt{(0.387)^3} = 0.240\,75\,\mathrm{yr} = 87.934\,\mathrm{days}
\]
\hfill(2.0)

We know that $\omega_{\mathrm{orb}} = v/r$ so will vary throughout the orbit (since $r$ varies within
an ellipse) whilst $\omega_{\mathrm{rot}}$ is constant.
\[
\omega_{\mathrm{rot}} = \frac{2\pi}{T_{\mathrm{rot}}}
= \frac{2\pi}{\frac{2}{3} \times 87.934 \times 86400}
= 1.24 \times 10^{-6}\,\mathrm{rad/s}
\]
\hfill(2.0)

For retrograde motion, we need $\omega_{\mathrm{orb}} \ge \omega_{\mathrm{rot}}$.\\
Using the vis-viva equation for the critical value of $r$ for when
$\omega_{\mathrm{rot}} = \omega_{\mathrm{orb}}$,
\begin{align*}
\omega_{\mathrm{orb}}^2 &= \frac{v^2}{r^2}
  = \frac{GM_\odot}{r^2}\left(\frac{2}{r} - \frac{1}{a}\right) = \omega_{\mathrm{rot}}^2\\
\frac{2GM_\odot}{r^3} - \frac{GM_\odot}{ar^2} &= \omega_{\mathrm{rot}}^2\\
\therefore \frac{\omega_{\mathrm{rot}}^2}{GM_\odot}r^3 + \frac{1}{a}r - 2 &= 0\\
(1.160\times10^{-32})r^3 + (1.725\times10^{-11})r - 2 &= 0 \quad
  \text{(if $r$ is in meters)}\\
38.99r^3 + 2.584r - 2 &= 0 \quad \text{(if $r$ is in au)}
\end{align*}
\hfill(4.0)

Solving the cubic equation (by any valid method) gives only one non-imaginary root.
One possible way would be to use iterations for this polynomial $f(r)$, using perihelion
distance as the starting guess.
\begin{center}
\begin{tabular}{c|c}
$r$ & $f(r)$\\
\hline
0.3073 & $-0.07482$\\
0.3100 & $-0.03747$\\
0.3120 & $-0.00967$\\
0.3140 & 0.01841\\
0.3130 & 0.00433\\
0.3127 & 0.00012\\
\end{tabular}
\end{center}
\[
\therefore r = 0.3127\,\mathrm{au} = 4.684\times10^{10}\,\mathrm{m}
\]
\hfill(4.0)

From knowledge of ellipses, if $E$ is the eccentric anomaly
\begin{align*}
r &= a(1 - e\cos E)\\
\therefore E &= \cos^{-1}\left[\frac{1}{e}\left(1 - \frac{r}{a}\right)\right]
  = \cos^{-1}\left[\frac{1}{0.206}\left(1 - \frac{0.3127}{0.387}\right)\right]\\
E &= 0.3706\,\mathrm{rad} = 21.23^\circ = 21^\circ 14'
\end{align*}
\hfill(3.0)

Using Kepler's Equation we can find the mean anomaly, $M$
\begin{align*}
M &= E - e\sin E = 0.3706 - 0.206 \times \sin 0.3706\\
  &= 0.2960\,\mathrm{rad} = 16.96^\circ = 16^\circ 58'
\end{align*}
\hfill(2.0)

This is relative to the perihelion, so symmetry demands that the total time the sun is in
retrograde corresponds to
\begin{align*}
\Delta M &= 2M\\
\therefore T_{\odot,retro} &= T_{\mathrm{Mercury}}\frac{\Delta M}{2\pi}
  = 87.934 \times \frac{2 \times 0.296}{2\pi}\\
T_{\odot,retro} &= \boxed{8.28\,\mathrm{days}}
\end{align*}
\hfill(3.0)

[Accept alternative methods making use of the true anomaly e.g.\ with the relation
$\cos E = \frac{e + \cos\nu}{1 + e\cos\nu}$.]
